\documentclass[conference]{IEEEtran}

\usepackage{amsmath, amssymb}
\usepackage{booktabs}
\usepackage{hyperref}
\usepackage{soul}

\hypersetup{
  colorlinks=true,
  linkcolor=black,
  urlcolor=black
}

\begin{document}

\title{A Sequential Decision-Making Framework for a Two-Player Resource-and-Combat Game\\
{\normalsize ASEN 5264 --- Decision Making Under Uncertainty}}

\author{
  \IEEEauthorblockN{Michael Sola}
  \IEEEauthorblockA{University of Colorado Boulder}
  \and
  \IEEEauthorblockN{Jeremiaha Hoffman}
  \IEEEauthorblockA{University of Colorado Boulder}
}

\maketitle

% -----------------------------------------------------------------------
\section{Introduction}
% -----------------------------------------------------------------------

This project develops a two-player strategy game in which each player manages workers,
marines, and resources with the goal of eliminating all enemy marines. The game is
formalized as a Markov Decision Process and solved using methods from the course, with
increasing model complexity across three tiers of completion. The first tier utilized
value iteration to solve for an optimal policy against a deterministic policy of the
adversary and reinforcement learning, specifically Q-learning, Deep Q-Network (DQN)
learning, and \hl{add others}, to learn against this optimal policy. The second tier
formulated this problem as both a Markov Game (MG) and a Partially Observable Markov
Decision Process (POMDP). The third tier combined the MG and POMDP into a Partially
Observable Stochastic Game (POSG). \hl{summarize results}

% -----------------------------------------------------------------------
\section{Background and Related Work}
% -----------------------------------------------------------------------

\hl{Provide at least 5 citations of related published work and briefly describe how each
relates to this project.}

% -----------------------------------------------------------------------
\section{Problem Formulation}
% -----------------------------------------------------------------------

\subsection{State Space}

The state for each player $i \in \{1, 2\}$ is $(w_i, m_i, r_i)$, representing workers,
marines, and resources. The joint state is $s = (w_1, m_1, r_1, w_2, m_2, r_2)$, plus
a terminal state $s_{\text{term}}$. Both players start at $w = m = r = 1$. Resources
increase by $w_i$ each turn.

\subsection{Action Space}

Each player selects from the following action set each turn:
\[
  \mathcal{A} = \{\texttt{train worker},\; \texttt{train marine},\; \texttt{attack}\}
\]

\subsection{Transition Model}

The table below summarizes the transition model. If either player has $m_i = 0$, the
game transitions to $s_{\text{term}}$ with probability 1 regardless of action.

\begin{table}[h]
\centering
\caption{Transition Model}
\begin{tabular}{llll}
  \toprule
  \textbf{Action} & \textbf{Condition} & \textbf{Next State} & \textbf{Prob.} \\
  \midrule
  train worker  & $r \geq 1$  & $w' = w+1,\; r' = r-1$                      & $0.9$ \\
  train marine  & $r \geq 1$  & $m' = m+r,\; r' = 0$                         & $0.9$ \\
  attack        & always      & $m'_A = m_A - k_A,\; m'_B = m_B - k_B$       & $P(k_A, k_B)$ \\
  terminal      & $m_i = 0$   & $s_{\text{term}}$                             & $1.0$ \\
  \bottomrule
\end{tabular}
\end{table}

where $k_A \sim \mathrm{Binomial}(m_B,\, p_B)$, $k_B \sim
\mathrm{Binomial}(m_A,\, p_A)$, and $P(k_A, k_B)$ is their product under independence.

\subsection{Reward}

$R = 1$ upon reaching $s_{\text{term}}$ as the winner; $R = 0$ otherwise.

% -----------------------------------------------------------------------
\section{Solution Approach}
% -----------------------------------------------------------------------

\subsection{Tier 1 --- MDP, Value Iteration, and Reinforcement Learning}

The game is formalized as an MDP and solved via value iteration to obtain an optimal
policy $\pi^*$. A reinforcement learning agent (e.g., Q-learning or DQN) is
then trained on the same environment, with the two approaches compared in terms of
policy quality and convergence.

\subsection{Tier 2 --- Markov Game and POMDP}

The model is extended to either a Partially Observable MDP, where each player observes
only a noisy estimate of the opponent's state, or a Markov Game, where both players act
simultaneously and transitions depend on the joint action.

\subsection{Tier 3 --- Partially Observable Stochastic Game}

The full formulation combines partial observability with simultaneous multi-agent
interaction. Each player maintains a belief over the joint state and reasons about
opponent behavior under uncertainty. Solution methods may include multi-agent
reinforcement learning or approximate dynamic programming techniques.

% -----------------------------------------------------------------------
\section{Results}
% -----------------------------------------------------------------------

\hl{Present results with plots and tables. Quantify uncertainty where possible (e.g.,
confidence intervals over multiple runs). Compare value iteration policy against RL
agents.}

% -----------------------------------------------------------------------
\section{Conclusion}
% -----------------------------------------------------------------------

\hl{Summarize findings, key takeaways, and potential future extensions.}

% -----------------------------------------------------------------------
\section*{Contributions}
% -----------------------------------------------------------------------

\hl{Describe the contributions of each team member. Example: ``Michael implemented value
iteration and the MDP environment; Jeremiaha implemented the RL agents and ran
experiments.''}

The authors grant permission for this report to be posted publicly.

\bibliographystyle{IEEEtran}
\bibliography{references}

\end{document}
